\chapter{Atenção ao Estudante}
\label{chp:atendimento_estudante}

O \ppccurso da UFU compreende que a excelência acadêmica exige um ambiente de aprendizagem acolhedor, saudável, diverso e acessível. Por isso, desenvolve e apoia ações de atenção estudantil em articulação com a Faculdade de Engenharia Elétrica (FEELT) e os órgãos centrais da universidade.

\section{Acompanhamento Acadêmico e Pedagógico}

A Coordenação do Curso e a Direção da FEELT acompanham ativamente o percurso dos(as) estudantes, prestando orientação individual, apoio à matrícula, aconselhamento pedagógico e mediação de conflitos acadêmicos. Também são promovidas ações de recepção de calouros e oficinas de ambientação universitária.

Para fortalecer esse suporte, a FEELT instituiu um núcleo específico para o acolhimento e a orientação da comunidade discente.

\section{Orientação Acadêmica na Área Básica de Ingresso (ABI)}
\label{sec:orientacao_academica_abi}

Se a proposta de ABI for aprovada nos termos do \autoref{chp:abi}, a Coordenação do Curso e a Direção da FEELT prestarão, ao longo dos cinco primeiros períodos, orientação acadêmica específica, incluindo:

\begin{itemize}
    \item o esclarecimento, desde o ingresso, sobre a existência das duas graduações integrantes da ABI, suas trajetórias específicas e seus respectivos perfis de egresso (\autoref{chp:perfil_egresso});
    \item o apoio à escolha da graduação a ser concluída inicialmente, a se formalizar ao final do quinto período (\autoref{sec:escolha});
    \item a informação clara sobre os limites da ABI e sobre eventual mecanismo de continuidade entre cursos, somente se este vier a ser regulamentado (\autoref{sec:permanencia}); e
    \item o acompanhamento de eventual continuidade autorizada, inclusive quanto ao procedimento formal de aproveitamento de estudos (\autoref{sec:equivalencias_abi}).
\end{itemize}

\section{Núcleo de Apoio e Atenção ao Estudante da FEELT (NAAE-FEELT)}

Em conformidade com a política institucional de assistência estudantil~\cite{ufu_consex_20_2022} e seu regimento interno~\cite{ufu_feelt_regimento_2023}, a Faculdade de Engenharia Elétrica dispõe do Núcleo de Apoio e Atenção ao Estudante (NAAE-FEELT). Este núcleo atua como um ponto de referência para o acolhimento, a escuta qualificada e o encaminhamento das demandas estudantis no âmbito da faculdade.

Os principais objetivos do NAAE-FEELT são:
\begin{itemize}
\item Propor e executar ações de acolhimento e integração para estudantes ingressantes e veteranos;
\item Oferecer suporte pedagógico e psicossocial em primeira instância, atuando de forma preventiva e mediadora;
\item Articular as demandas dos estudantes com as coordenações de curso e os serviços especializados da UFU, como a Pró-Reitoria de Assistência Estudantil (PROAE);
\item Contribuir para a redução dos índices de retenção e evasão, promovendo a permanência qualificada;
\item Fomentar um ambiente acadêmico inclusivo, plural e respeitoso.
\end{itemize}

As ações a serem desenvolvidas incluem recepção de calouros, rodas de conversa, oficinas sobre métodos de estudo e saúde mental, além de atendimentos individuais e coletivos.

Mais informações:

{\footnotesize\url{https://www.feelt.ufu.br/unidades/nucleo/nucleo-de-apoio-e-atencao-ao-estudante}}

\section{Assistência Estudantil e Políticas de Permanência}

A Pró-Reitoria de Assistência Estudantil (PROAE) oferece suporte amplo aos(às) estudantes, com ações conduzidas por diretorias especializadas como a DIRES (Diretoria de Inclusão, Promoção e Assistência Estudantil) e DIPAE (Divisão de Promoção de Igualdade e Apoio Educacional). Entre os programas e serviços, destacam-se:

\begin{itemize}
\item Auxílios financeiros para moradia, alimentação, transporte e creche;
\item Bolsas acadêmicas (monitoria, PIBIC, extensão, PET);
\item Apoio à acessibilidade e à inclusão de pessoas com deficiência;
\item Acompanhamento psicopedagógico e social individualizado;
\item Incentivo à inclusão digital e ao acesso universal ao conhecimento;
\item Reconhecimento da diversidade étnico-racial, cultural e de gênero.
\end{itemize}

\section{Saúde Mental e Bem-Estar Psicológico}

A UFU mantém serviços de orientação em saúde mental por meio da Divisão de Saúde do Estudante (DISAU), que atua em parceria com o Sistema Único de Saúde (SUS) e projetos de promoção de bem-estar. O atendimento contempla:

\begin{itemize}
\item Apoio psicológico individual e escuta qualificada;
\item Encaminhamento para a rede SUS quando necessário;
\item Campanhas educativas e preventivas (ex: ansiedade, depressão, suicídio);
\item Oficinas temáticas e grupos de apoio organizados pelo setor especializado.
\end{itemize}

Mais informações:

{\footnotesize\url{https://proae.ufu.br/servicos/orientacao-em-saude-mental}}

\section{Promoção da Diversidade, Inclusão e Direitos Humanos}

A UFU reafirma seu compromisso com o combate a todas as formas de preconceito, discriminação e exclusão. Campanhas institucionais e programas permanentes promovem o respeito às diferenças e a construção de um ambiente universitário mais justo, seguro e plural.

Ações voltadas ao combate à LGBTfobia, ao racismo, ao capacitismo e à xenofobia são desenvolvidas em articulação com o Comitê de Diversidade da UFU, os coletivos estudantis e os programas institucionais da PROAE, DIEPAFRO e CPDiversa.

\section{Canais de Apoio ao Estudante}

A Universidade Federal de Uberlândia oferece diversos serviços de apoio acadêmico, psicossocial e institucional para garantir a permanência qualificada dos(as) estudantes. Abaixo, estão relacionados alguns dos principais canais de apoio:

\begin{longtblr}[
theme = ppc,
caption = {Principais Canais de Apoio ao Estudante},
label = {tab:canais_apoio_estudante},
]{
colspec = {Q[l,m,wd=30mm]Q[c,m,wd=35mm]Q[c,m,wd=70mm]},
rowhead = 1,
row{odd} = {bg=CinzaClaro},
row{1} = {bg=AzulEscuro, fg=white},
cells  = {font = \fontsize{10pt}{12pt}\selectfont},
}
\textbf{Área de Apoio} & \textbf{Serviço / Portal} & \textbf{Link / Contato}\\
Apoio Local (FEELT) & Núcleo de Apoio e Atenção ao Estudante (NAAE-FEELT) & \url{https://www.feelt.ufu.br/unidades/nucleo/nucleo-de-apoio-e-atencao-ao-estudante}\\
% REVISAR: página institucional própria deste curso (curso em fase de proposta) não localizada — confirmar URL com a Coordenação antes da submissão deste PPC.
Suporte à vida acadêmica & Coordenação do Curso de \ppccurso (FEELT) & \url{https://www.feelt.ufu.br}\\
Apoio pedagógico individual e coletivo & Divisão de Promoção de Igualdades e Apoio Educacional (DIPAE/PROAE) & \url{https://proae.ufu.br/dipae}\\
Apoio psicossocial e saúde mental & Divisão de Saúde do Estudante (DISAU/PROAE) & \url{https://proae.ufu.br/servicos/orientacao-em-saude-mental}\\
Inclusão, permanência e assistência estudantil & Diretoria de Inclusão, Promoção e Assistência Estudantil (DIRES) e Divisão de Promoção de Igualdade e Apoio Educacional (DIPAE) & \url{https://proae.ufu.br/dipae}\\
Ações de respeito e diversidade & DIEPAFRO, CPDiversa, Campanhas e ações institucionais & \url{https://comunica.ufu.br}\\
Sistema de Saúde Universitário & Divisão de Saúde do Estudante (DISAU/PROAE) & \url{https://proae.ufu.br/disau}\\
Mobilidade acadêmica nacional & Programa ANDIFES & \url{https://prograd.ufu.br/mobilidadenacional}\\
Mobilidade acadêmica internacional & Diretoria de Relações Internacionais e Interinstitucionais (DRII) & \url{https://dri.ufu.br/}\\
Sistemas e documentos acadêmicos & Portal do Estudante de Graduação & \url{https://www.portalestudante.ufu.br/}\\
Portal de notícias & Informações institucionais e notícias & \url{https://comunica.ufu.br/}\\
Ouvidoria e denúncias & Canal de Ouvidoria & \url{https://ufu.br/ouvidoria-geral}\\
\end{longtblr}

A UFU conta com comissões permanentes, grupos estudantis, projetos de extensão e programas de mentoria voluntária que fortalecem a escuta e o acolhimento à comunidade discente. É fundamental que o(a) estudante acompanhe os canais oficiais da PROAE e da FEELT para informações sobre editais e cronogramas.
